\section*{第 \ref{ch:b1:charged-particles} 章——电场与磁场中的带电粒子}

\begin{solution}{exo:b1:charged-particles:1}
$\qty{1}{eV} = \qty{1.60e-19}{J}$, $\qty{13.6}{eV} = \qty{2.18e-18}{J}$,
$\qty{1}{MeV} = \qty{1.60e-13}{J}$。$v = \sqrt{2eU/m}$：电子
\qty{1.9e7}{m/s}；质子 \qty{4.4e5}{m/s}（慢 $\sqrt{1836}$ 倍）。
\end{solution}

\begin{solution}{exo:b1:charged-particles:2}
$v = \qty{1.88e7}{m/s}$; $R = mv/eB = 9.11\times10^{-31} \times 1.88\times10^7/
(1.6\times10^{-19} \times 10^{-3}) = \qty{0.11}{m}$; $\omega_c = eB/m =
\qty{1.8e8}{rad/s}$; $T = 2\pi/\omega_c = \qty{36}{ns}$.
\end{solution}

\begin{solution}{exo:b1:charged-particles:3}
$R = mv/qB = 1.67\times10^{-27} \times 10^5/(1.6\times10^{-19} \times 5\times
10^{-5}) = \qty{21}{m}$。在 $30^\circ$ 时，半径为 $R\sin 30^\circ
= \qty{10}{m}$ 的螺旋线，螺距为 $2\pi(m/qB)v\cos 30^\circ = \qty{114}{m}$，
沿磁场线漂移。
\end{solution}

\begin{solution}{exo:b1:charged-particles:4}
$v = E/B = 10^4/0.020 = \qty{5.0e5}{m/s}$。速率较大的离子：磁力
占优，向磁力一侧偏转；电荷相反而
速率相同的粒子：两力同时反向，仍然平衡——选择器
对电荷符号和质量都不敏感。
\end{solution}

\begin{solution}{exo:b1:charged-particles:5}
$\tan\alpha = EL/2U_0 = 10^4 \times 0.04/4000 = 0.10$；偏移量
$0.30 \times 0.10 = \qty{3.0}{cm}$；灵敏度 \qty{0.30}{mm/V}。较快的
电子在极板间停留时间较短：$\tan\alpha \propto 1/U_0$。
\end{solution}

\begin{solution}{exo:b1:charged-particles:6}
$R = \sqrt{2mU/e}/B$: $m_{12} = \qty{1.99e-26}{kg}$, $R_{12} = \sqrt{2 \times
1.99\times10^{-26} \times 2\times10^4/1.6\times10^{-19}}/0.5 = \qty{14.1}{cm}$;
$R_{13} = R_{12}\sqrt{13/12} = \qty{14.7}{cm}$；落点间距 $2\Delta R =
\qty{1.2}{cm}$（两落点的半径相差 \qty{5.8}{mm}，
直径相差其两倍）。
\end{solution}

\begin{solution}{exo:b1:charged-particles:7}
$f_c = eB/2\pi m = \qty{23}{MHz}$; $E_k = e^2B^2R^2/2m = \qty{4.3e-12}{J} =
\qty{27}{MeV}$；每圈 \qty{100}{keV}：$270$ 圈；时间 $270/f_c =
\qty{12}{\micro s}$; $\gamma - 1 = 27/938 = 2.9\%$.
\end{solution}

\begin{solution}{exo:b1:charged-particles:8}
$v_\perp = \qty{5.0e5}{m/s}$, $v_\parallel = \qty{8.7e5}{m/s}$: $R = mv_\perp/eB
= 9.11\times10^{-31} \times 5\times10^5/(3.2\times10^{-22}) = \qty{1.4}{mm}$;
周期 $2\pi m/eB = \qty{18}{ns}$，螺距 $v_\parallel T = \qty{1.5}{cm}$：
每米约 $64$ 圈。
\end{solution}

\begin{solution}{exo:b1:charged-particles:9}
$U_H = IB/nqd = 5.0/(8.5\times10^{28} \times 1.6\times10^{-19} \times 10^{-3})
= \qty{3.7e-7}{V}$：亚微伏量级。若 $n$ 小千万倍，
则半导体可得毫伏量级\omterm{def:b1:dc-circuits:voltage}{电压}——一种实用的传感器。
\end{solution}

\begin{solution}{exo:b1:charged-particles:10}
$m\,\dd\vect v/\dd t = q\vect v\wedge\vect B$；与 $\vect v$ 点乘：$\dd(\tfrac12 mv^2)/
\dd t = 0$，无论 $\vect B(\vect r, t)$ 如何，只要没有 $\vect E$。
速度的方向、\omterm{def:b1:point-kinematics:frenet}{曲率半径}、螺距都可能
改变；速率和\omterm{thm:b1:work-and-energy:ke}{动能}（\qty{10}{keV}）则不能。
\end{solution}

\begin{solution}{exo:b1:charged-particles:11}
$m\dot v_x = qv_yB$，$m\dot v_y = qE - qv_xB$，$v_z = 0$。取 $\vect v =
(E/B)\vect e_x + \vect w$：$m\dot w_x = qw_yB$，$m\dot w_y = -qw_xB$ ——
即 $\vect w$ 以 $\omega_c$ 做纯粹的\omterm{thm:b1:charged-particles:bfield}{回旋运动}。从静止出发，
$\vect w(0) = -(E/B)\vect e_x$：$\vect w$ 的模恒为 $E/B$，于是
速度等于漂移速度 $E/B$ 与一个模相同的旋转矢量之和——
摆线，在 $\vect v = \vect 0$ 处出现尖点。电子：漂移速度 $E/B =
\qty{1.0e5}{m/s}$，$\omega_c = \qty{1.8e10}{rad/s}$，拱高 $2mE/qB^2 =
\qty{11}{\micro m}$.
\end{solution}

\begin{solution}{exo:b1:charged-particles:12}
$R = mv_0/eB = 9.11\times10^{-31} \times 2.65\times10^7/(1.6\times10^{-22}) =
\qty{0.15}{m}$；$\alpha \approx L/R = 0.04/0.15 = 0.27$（约 $15^\circ$）：
偏移量 $\approx 0.30 \times 0.27 = \qty{8}{cm}$，所需磁场仅为适中的
毫特斯拉量级。若要用极板获得大偏转角，则需千伏\omterm{def:b1:dc-circuits:voltage}{电压}和较长的
极板；而线圈可使束流偏转数十度，
且完全不需要\omterm{def:b1:dc-circuits:voltage}{电压}。
\end{solution}

\begin{solution}{pb:b1:charged-particles:1}
\textbf{1.} $\vect F = e\vect v\wedge\vect B \perp \vect v$：不做功，速率
恒定。

\textbf{2.} $mv^2/R = evB$：$R = mv/eB$；$T = 2\pi R/v = 2\pi m/eB$，与
$v$ 无关。

\textbf{3.} $f_c = eB/2\pi m = 1.6\times10^{-19} \times 1.5/(2\pi \times 1.67
\times10^{-27}) = \qty{23}{MHz}$.

\textbf{4.} $eU = \qty{50}{keV}$ 每次穿越间隙，\qty{100}{keV} 每圈。

\textbf{5.} 每次质子到达时，间隙都必须处于加速极性，
两次到达相隔半个周期；频率偏离时，\omterm{def:b1:wave-propagation:signal}{相位}会滑移，
最终的加速冲量会使粒子减速。

\textbf{6.} 绕行一圈后 $E_k = \qty{100}{keV}$，$v = \sqrt{2E_k/m} =
\qty{4.4e6}{m/s}$，$R = mv/eB = \qty{3.0}{cm}$；绕行 $n$ 圈后 $E_k =
n \times \qty{100}{keV}$，$R = \qty{3.0}{cm}\sqrt n$。

\textbf{7.} $E_k = e^2B^2R^2/2m = \qty{4.3e-12}{J} = \qty{27}{MeV}$; $v =
\sqrt{2E_k/m} = \qty{7.2e7}{m/s} = 0.24c$.

\textbf{8.} 共 $270$ 圈；$270 \times \qty{43.5}{ns} = \qty{12}{\micro s}$。

\textbf{9.} $\sum 2\pi R_n = 2\pi \times 0.030\sum\sqrt n \approx 2\pi \times 0.030
\times \tfrac23 \times 270^{3/2} = \qty{560}{m}$.

\textbf{10.} 否：$E_k$ 仅取决于 $B$ 与 $R_{\max}$；$U$ 决定
圈数与时间。

\textbf{11.} 路径长 \qty{560}{m}：需 $\lambda \gtrsim \qty{1}{km}$，即\
$p \lesssim 10^5\ \unit{Pa} \times 7\times10^{-8}/10^3 = \qty{7e-6}{Pa}$ ---
高真空。

\textbf{12.} $\gamma - 1 = 27/938 = 0.029$；周期增加 $2.9\%$。
经过 $270$ 圈后，滑移累积到数个周期——末尾几圈
将失去同相；实际中 $B$ 与注入\omterm{def:b1:wave-propagation:signal}{相位}
都需微调，\qty{27}{MeV} 已接近该简单机械的极限。

\textbf{13.} $E_k = (eBR)^2/2m$: $BR = \sqrt{2mE_k}/e = \sqrt{2 \times 1.67\times
10^{-27} \times 3.68\times10^{-11}}/1.6\times10^{-19} = \qty{2.2}{T.m}$; $R =
\qty{1.1}{m}$，在 \qty{2}{T} 下。

\textbf{14.} $T = 2\pi\gamma m/eB$ 由中心到边缘增加 $24\%$：一
个固定频率将彻底失步。

\textbf{15.} 同步回旋加速器：束团获得能量时降低频率
——任一时刻仅一个束团，平均电流小。等时性回旋加速器：
将磁场设计为 $B(R) \propto \gamma$，使周期保持恒定——连续
束流，是现代的选择。

\textbf{16.} $p = \gamma mv = 1.245 \times 1.67\times10^{-27} \times 1.8\times10^8
= \qty{3.7e-19}{kg.m/s}$；$R = p/eB = \qty{1.2}{m}$，略大于
非相对论的 \qty{1.1}{m}。

\textbf{17.} 横向冲量 $eEL/v = 1.6\times10^{-19} \times 10^5 \times
1.0/1.8\times10^8 = \qty{8.9e-23}{kg.m/s}$；角度 $\approx 8.9\times10^{-23}/
3.7\times10^{-19} = \qty{2.4e-4}{rad}$：转向毫无希望。

\textbf{18.} 磁偏转：$R = p/eB = \qty{2.3}{m}$，角度 $L/R = \qty{0.43}{rad}
\approx 25^\circ$：比电偏转大一千倍——故采用磁铁。

\textbf{19.} 每个质子的能量 $\qty{3.68e-11}{J}$；剂量 $2.0 \times 0.50 =
\qty{1.0}{J}$：约需 $2.7\times10^{10}$ 个质子。

\textbf{20.} $10^{-9}/1.6\times10^{-19} = \qty{6.2e9}{protons/s}$：约
\qty{4.4}{s}。

\textbf{21.} $P = \qty{230}{MeV} \times \qty{6.2e9}{s^{-1}} = \qty{0.23}{W}$:
夜灯，恰好输送到所需之处。

\textbf{22.} $\qty{58}{MeV}$ 在 \qty{1}{cm} 内，约 \qty{6}{keV/\micro m}——
沿末段径迹电离密集。

\textbf{23.} 每个质子约产生 $230\times10^6/30 \approx \num{8e6}$ 次电离——
造成细胞死亡的分子损伤。

\textbf{24.} X 射线沿路径各处都沉积能量，其中多数集中在
入射端附近；质子则把大部分能量沉积在射程末端，
从而保护前方组织，射程之后不留沉积。

\textbf{25.} $R = mv/eB$ 与 $T = 2\pi m/eB$（圆周），$E_k = (eBR)^2/2m$
（能量），$E_k/2eU$（圈数），剂量 $=$（质子数 $\times$ 能量）/质量
以及电流 $= e \times$ 速率（时间）。
\end{solution}
